Erklären Sie die initialen Funktionen von Versicherungen im Rahmen von Versicherungsrecht

1. Sicherungsfunktion (vgl. § 1 VVG)
Versicherung = Instrument des Risikomanagement è
Risikotransfer statt Risikoverminderung („weniger handeln“) oder
Risikovermeidung („gar nicht handeln“)

2. Innovationsfunktion
Versicherung erlaubt Zukunftsplanung, Handeln, Wirtschaften:
• Im privaten Bereich: z.B. Hausbau (Lebens-/RestschuldVers)
• Im unternehmerischen Bereich: (Medizin)Produktherstellung
(ProdukthaftpflichtVers)
• Entwicklung weltweiten Handels und modernen Industrie setzt
Versicherungswesens voraus (Wurzeln: insbes. im Seehandel)
• Daher: Hochgradig relevant, hochgradig international


Von welchen Versicherungen ist das Privatversicherungsrecht abzugrenzen?

Sozialversicherungen!
• Versicherungsverhältnis entsteht nicht durch Vertragsschluss,
sondern durch Gesetz (z.B. „Mitgliedschaft“ in gesetzlicher
Renten- und Krankenversicherung sowie Pflegeversicherung)
• Grundsätzlich auch Umlageverfahren und Beitragssystem (z.B. bei
der gesetzlichen Krankenversicherung, wobei hier „Substitutive“-
PKV möglich, vgl. § 195 VVG)


Wie lautet die Gliederung des Privatversicherungsrecht?

Privatversicherungsrecht als Materie umfasst:
• Versicherungsvertragsrecht (VVG): 
Regelt Rechtsbeziehungen zw. Versicherer (VR) und
Versicherungsnehmer (VN)/Versicherten; zudem Recht der
Versicherungsvermittlung zw. Vermittler und VN

• Versicherungsaufsichtsrecht (insbes. VAG): 
Öffentl. R und befasst sich mit Zulassung/Kontrolle der VR

• Versicherungsunternehmensrecht: Gründung, Gestaltung
und Umstrukturierung von VR 
(VAG, Gesellschaftsrecht und teils Vereinsrecht bei VVaG)


Kategorisieren Sie aus rechtlicher Sicht die Arten von Versicherungsveträgen

Die Systematik des VVG orientiert sich an der Unterscheidung nach Schadens-, Summen-, Personen- und Sachversicherung. 
Die §§ 74 ff. VVG beziehen sich nur auf die Schadensversicherung, wobei §§ 88 ff. VVG nur die Schadensversicherung als Sachversicherung regeln. 
Die Kapitel 5-8 beziehen sich nur auf die Personenversicherung. 
Bei den Versicherungsverträgen wird grundsätzlich zwischen Schadens- und Summenversicherungen unterschieden.

Beispiele: 
Schadensversicherungen: 
Haftpflichtversicherung, 
Hausratversicherung, 
Feuerversicherung, 
Summenversicherung:
Unfallver.,
Berufsunfähigkeitver.,
Lebensver.

Die Schadensversicherung wird in eingeteilt in
die Aktivversicherung (Schutz einzelner Vermögensbestandteile; spezifisches Schutzobjekt) 
oder auch beschrieben als 
Ausgleich des entstandenen Verlusts (Beschädigung) von Vermögenswerten.

und in die Passivenversicherung (Schutz des gesamten Vermögens)
oder auch beschrieben als Ausgleich von entstandenen Verbindlichkeiten.

Beispiele: 
Aktivversicherung: 
Forderungsversicherung, Sachversicherung;
Passivenversicherung: 
Haftpflichtversicherung, Rechtsschutzversicherung

Ferner differenziert man zwischen Personenversicherungen (z.B. Lebensversicherung - Versicherung gegen Todesrisiko einer Person) und 
Sachversicherungen (z.B. Hausratsversicherung - Sicherung des Hausrates gegen Einbruchdiebstahl). 
Der Unterschied besteht im Risikoträger bzw. Risikoobjekt.

Im Rahmen der Kategorisierung ist noch die Pflichtversicherung zu nennen. 
In der Kfz-Haftpflichtversicherung (vgl. § 1 PflVG) und der Krankenversicherung zum Basistarif (vgl. § 193 V 1 VVG) trifft den Versicherer 
ein Kontrahierungszwang.


Was sind die 3 Komponenten des Versicherungsgeschäfts

Risikogeschäft (Gewähr einer Versicherungsleistung und die Abgabe des Versicherungsschutzversprechens)
Spar- und Enstpargeschäft (e.g. kapitalbildende Lebensversicherung)
Dienstleistungsgeschäft (Abwicklung des Versicherungsfalles, Hotline)


Was sind die Merkmale von versicherbaren Risiken?

Zufälligkeit des Schadeneintritts
Schätzbarkeit der Schadensverteilung
Eindeutige Definition des Risikos möglich
Beherrschbarer Höchstschaden
Unabhängigkeit der Einzelrisiken/kein Kumulschaden zur Anwendung des Gesetzes der großen Zahlen.


Was sind die Merkmale eines Versicherungsvertrages nach §1 VVG und die Konkretisierungs des BGH.

§ 1 VVG [Vertragstypische Pflichten]: Der Versicherer
verpflichtet sich mit dem Versicherungsvertrag, ein
bestimmtes Risiko des Versicherungsnehmers oder
eines Dritten durch eine Leistung abzusichern, die er bei
Eintritt des vereinbarten Versicherungsfalls zu erbringen
hat. Der Versicherungsnehmer ist verpflichtet, an den
Versicherer die vereinbarte Zahlung (Prämie) zu leisten

Laut BGH NJW 2017,
1. ungewisses Ereignis (Risiko: Gefahr, Zufälligkeit)
• zukünftiges oder ggfs. auch ein schon vergangenes Ereignis (§ 2 VVG)
• ungewiss ob oder wann („timing risk“) ein Ereignis eintritt
• = Risikoelement, aleatorisches Element
• 2. Leistungspflicht des VR
• Damit korrespondieren: Rechtsanspruch des VN auf Lstg im Versicherungsfall
• 3. Entgeltlichkeit
• Pflicht des VN zur Entrichtung von Prämie/Beitrag (§§ 33 ff VVG)
• 4. Planmäßigkeit
• Ausgleich gleichartiger Risiken im Kollektiv (Streuung in Risikogemeinschaft)
• durch Kalkulation der Prämie nach dem Gesetz der großen Zahl
• 5. Selbständigkeit der Risikotragung
• Keine Versicherungsgeschäft, wenn die Risikoübernahme in einem inneren
Zusammenhang mit einem anderen Rechtsgeschäft steht (zB kein Annexvertrag e.g. Kaufvertrag)

Geben Sie Beispiele von Geschäften die nicht Versicherungsveträge sein können
1. Derivate
2. Spiel und Wette
§ 762 BGB [Spiel, Wette]: „Durch Spiel oder durch
Wette wird eine Verbindlichkeit nicht begründet. Das auf
Grund des Spiels oder der Wette Geleistete kann nicht
deshalb zurückgefordert werden, weil eine verbindliche
Vereinbarung nicht bestanden hat.
3. Leibrente (Überlebensversicherung/Rente einer einzelnen Person)
§ 99 Gesetz über den Wertpapierhandel (WpHG): Ausschluss des
Einwands nach § 762 des Bürgerlichen Gesetzbuchs
Gegen Ansprüche aus Finanztermingeschäften, bei denen
mindestens ein Vertragsteil ein Unternehmen ist, [...]
kann der Einwand des § 762 des Bürgerlichen Gesetzbuchs nicht
erhoben werden. Finanztermingeschäfte im Sinne des Satzes 1
und der §§ 100 und 101 sind die derivativen Geschäfte im Sinne
des § 2 Absatz 3 und Optionsscheine.


Nennen Sie die relevantesten Paragraphen des BGB's für das VVG

1. Buch (Allg. Teil)
§ 123 Anfechtbarkeit wegen Täuschung oder Drohung
§§ 130 Wirksamwerden der Willenserklärung gegenüber Abwesenden
§ 150 II BGB: Änderungen oder Erweiterungen gelten als neuer Antrag

2. Buch Recht der Schuldverhältnise
§ 242 BGB
Leistung nach Treu und Glauben
Der Schuldner ist verpflichtet, die Leistung so zu bewirken, wie Treu und Glauben mit Rücksicht auf die Verkehrssitte es erfordern.

3. Buch (Sachenrecht)
ist insgesamt relevant für VVG §95 (nicht BGB). Im Speziellen
§ 305 Einbeziehung Allgemeiner Geschäftsbedingungen in den Vertrag
Allgemeine Geschäftsbedingungen sind alle für eine Vielzahl von Verträgen vorformulierten Vertragsbedingungen.

4. Buch Familienrecht
§§ 1372 ff. BGB Zugewinnausgleich in anderen Fällen
§§ 1373 Zugewinn ist der Betrag, um den das Endvermögen eines Ehegatten das Anfangsvermögen übersteigt.
§ 2 II Nr. 3 VersAusglG

5. Buch Erbrecht
§§ 159 VVG
§328 BGB Vertrag zugunsten Dritter


Erläutern Sie inwiefern die AVB die Rechtssituation im Detail klärt und ihren Aufbau
(2) Allgemeine Versicherungbedingungen (AVB)
- AVB = besondere Allgemeine Geschäftsbedingungen (AGB), die durch VR
gestellt werden
- Gliedern sich in Allgemeine + Besondere Versicherungsbed.
+Risikobeschreibungen (insb. in Haftpflichtversicherung).

Die notwendigen Information werden durch §§ 1-4 VVG-InfoV geregelt.
Nach § 10 VAG a.F. mussten AVB bis 2016 z.B. mindestens enthalten:
• Ereignisse, bei deren Eintritt der VR zur Leistung verpflichtet ist und die Fälle,
in denen die Leistungspflicht ausgeschlossen sein soll,
• Art, den Umfang und die Fälligkeit der Leistungen des VR
• Fälligkeit der Prämie und die Rechtsfolgen des Verzugs,
• vertragliche Gestaltungsrechte des VN und des VR, Obliegenheiten und
Anzeigepflichten vor und nach Eintritt des Versicherungsfalls,
• Anspruchsverlust im Falle der Versäumung von Fristen,
• Inländische Gerichtsstände,
• Grundsätze und Maßstäbe, wonach die Versicherten an den Überschüssen
teilnehmen.

Für deren Die Formulierung gilt das Leitbild: 
„wie ein durchschnittlicher VN ohne versicherungsrechtliche
Spezialkenntnisse bei verständiger Würdigung und aufmerksamer
Durchsicht und Berücksichtigung des erkennbaren
Sinnzusammenhangs (die Klausel) verstehen muss.“
Zweifel bei der Auslegung gehen zulasten des VR § 305c II BGB:
Branchenbezogene Fachkenntnisse werden dennoch berücksichtigt 
(Beispiel Gastronomie, COVID vs. Rinderwahnsinn).


Stellen sie zwingende und halbzwingende Normen im VVG gegenüber

(a) Zwingend sind Normen, von denen keine Abweichung
zulässig ist
- Beispiel: § 5 Abs. 4 VVG: Kein Verzicht auf AnfechtungsR

(b) Halbzwingend sind dabei die Normen, die gerade den
VN schützen. Erforderlich ist, dass Abweichung dem VN
nachteilig ist. Abweichung zum Nachteil des VR ist
zulässig
- VVG kennzeichnet solche Normen unter der Überschrift
„Abweichende Vereinbarungen“ in den §§ 18, 32, 67, 87,
112, 129, 171, 175, 191, 208 VVG.
Diese Normen sind Verbotsvorschriften iSd § 134 BGB ,
d.h. ein Verstoß gegen sie führt zur Nichtigkeit der
entgegenstehenden Regelung im VersV.


Erläutern Sie den Ablauf des Vertragsschlusses im VVG

A. Vertragsschluss
Der Abschluss des Versicherungsvertrages richtet sich grundsätzlich nach den allgemeinen Regeln des BGB. 
Das VVG enthält jedoch einige Sonderregeln, wie die Beratungs- und Informationspflichten des Versicherers nach §§ 6, 7 VVG 
oder das Widerrufsrecht gemäß §§ 8, 9 VVG, welche den allgemeinen BGB-Vorschriften als lex specialis vorgehen.

I. Qualifiziertes Antragsmodell
Versicherungsvertrag, zwei korrespondierende Willenserklärungen, Antrag und Annahme, nach § 145 ff. BGB

Das VVG geht in § 7 von einem qualifizierten Antragsmodell aus. 
Der Gesetzgeber betrachtet somit das Antragsmodell als Regelfall des Vertragsschlusses. 
Foglich Antrag auf Abschluss des Versicherungsver. durch Versicherungsnehmer, nachdem zuvor der Versicherer diesen umfassend 
informiert hat. 
Gemäß § 7 I 1 VVG, Versicherer muss dem Versicherungsnehmer rechtzeitig vor Abgabe Vertragserklärung, Vertragsbestimmungen einschließlich 
der AVB sowie VVG-InfoV bestimmten Informationen (z.B. ladungsfähige Anschrift des Versicherers, Handelsregistereintrag) in Textform mitzuteilen. 
Die Definition von rechtzeitig ist nicht konkretisiert, der formularmäßige Informationsverzicht (§307 BGB) ist umstritten.

Wird der Versicherungsvertrag nicht unter ausschließlicher Verwendung von Fernkommunikationsmitteln abgeschlossen, 
kann der Versicherungsnehmer vor Abgabe der Vertragserklärung weitere Informationen verlangen, die der Versicherer ihm ebenfalls in 
Textform mitteilen muss (§ 7 I 2 VVG).

Die Annahme erfolgt durch den Versicherer, sie kann ausdrücklich oder konkludent (stillschweigende Willenserklärung, 
e.g. Kaufvertrag durch Ware auf Kasse legen) erfolgen.
Eine verspätete Annahme/Übersendung des Versicherungsscheines ist ein neuer Antrag seitens des Versicherers

II. Invitationmodell
Invitatio ad offerendum, Versicherungsnehmer fordert Versicherer zur Abgabe eines Angebotes auf.
Damit wird Versicherer zu Antragsteller.
Das Schweigen oder Duldung der ersten Prämie stellt keine konkludente Annahme des Versicherungsnehmer dar.

III. Policenmodell
Gängige Methode vor 2007, abgeschaft durch §7 VVG.

Die Details zum elektronischen Geschäftsverkehr regelt §312i BGB


Was gilt bei Abweichungen zw. Angebot (Antrag) und Annahme nach BGB?

§ 150 II BGB: Änderungen oder Erweiterungen gelten als neuer Antrag!
VVG sieht dagegen eine „Billigungsklausel“ in § 5 VVG als lex specialis zu § 150 II BGB vor:
§ 5 I VVG Abweichung „gilt als genehmigt“, wenn zusätzl Vorrs des Abs. 2 erfüllt (= Hinweis VR)
Sonst: § 5 III VVG: Inhalt des Antrags maßgeblich!
§ 5 IV VVG: VN kann wg Irrtums anfechten


Was ist das Schema für das Widerrufsrecht
Widerrufsrecht (§ 8 VVG)
1. Kein Ausschluss des Widerrufsrechts nach § 8 III VVG (z.B. bei Versicherungsverträgen über ein Großrisiko im Sinn des § 210 Absatz 2.)
2. Ordnungsgemäße Widerrufserklärung des VN
   - Auslegung des Willens lässt auf Vertragslösung schließen
   - in Textform, auch per E-Mail möglich
   - keine Begründung erforderlich
3. Frist
   - grundsätzlich zwei Wochen, § 8 I 1 VVG (30 Tage bei Lebensversicherung)
   - Absendung des Widerrufes genügt zur Fristwahrung

Die Rückerstattung von Prämien ist üblicherweise awendbar auf jene die nach Widerruf gezahlt wurden.



Was sind die Möglichkeiten zur Beendigung eines Versicherungsvertrages?

Es gelten die allg. Varianten nach BGB für Dauerschuldverhältnisse: 
Zeitablauf, Vertragsaufhebung, Kündigung, Rücktritt.

Ordentliche Kündigung § 11 VVG

Für automatische Vertragsverlängerungen gilt: 
Befristeter Vertrag, keine Verlängerung > 1 Jahr § 11 I VVG
Bei Vertragsdauer > 3 Jahre: 3 Monate zum Ende des 3. Jahres

Unbefristeter Vertrag: gem. § 11 III VVG nicht weniger als 1 und nicht mehr
als 3 Monate Kündigungsfrist

Außerordentliche Kündigung
In besonderen Fällen vorgesehen, z.B. § 40 I 1 VVG bei Prämienerhöhung
wg Anpassungsklausel, ohne dass sich der Versicherungsschutz
entsprechend ändert

Fristlose Kündigung bei besonders schweren Vertragsstörungen, wie zB
vorsätzliche oder grob fahrlässige Gefahrerhöhung durch VN: §24 I 1 VVG

Kündigungsrecht des VR
d) Rücktritt
• Bei vertragl Rücktrittsvorbehalt oder gesetzl Rücktrittsgrund
• Gesetzl RücktrittsR des VR:
• Bei nicht (rechtzeitiger) Erst-Prämien-Zahlung durch VN, § 37 I VVG.
Rechtsfolge neben § 346 BGB: Nach § 39 I 3 VVG behält VR angemessene
Geschäftsgebühr
• Wegfall des versicherten Interesses
• Bei Verletzung der vorvertragl Anzeigepflicht, §§ 19 I, 31 VVG.
Rechtsfolge neben § 346 BGB: VR steht nach § 39 I 2 VVG Prämie bis zum
Wirksamwerden des Rücktritts zu



